% Work/MAT — 1. ročník — hodina 60: Definiční obor a obor hodnot
% Návaznost: 59_FNC_CoJeFunkce; dále 61_FNC_GrafFunkceAPruseciky.
% Učebnice: ucebnice/02_Funkce_a_rovnice/Zaklady_funkci/Funkce.tex,
%   kapitola „Definiční obor a obor hodnot“; návaznost na lomené výrazy
%   v ucebnice/01_Zakladni_poznatky/Algebraicke_vyrazy/.
% Metodické podklady: Realisticky.cz, 020104 Funkce, definiční obor;
%   020105 Graf funkce I (důraz na čtení obou množin z grafu).
% Výklad, výpočty, úlohy i grafy jsou vlastní tvorba.
% Revize: přehledová tabulka podmínek včetně trigonometrických funkcí.
\documentclass[11pt,a4paper]{article}
% Spolecny styl pro ucitelske pripravy MAT (Work/MAT)
% Verze 1.1, 2026-10-23
\NeedsTeXFormat{LaTeX2e}
\ProvidesPackage{priprava}[2026/10/23 v1.1 Sablona ucitelskych priprav MAT]

% Zakladni typografie a cestina
\RequirePackage[T1]{fontenc}
\RequirePackage[utf8]{inputenc}
\RequirePackage[czech]{babel}
\RequirePackage{lmodern}
\RequirePackage{microtype}
\RequirePackage{amsmath,amssymb}
\RequirePackage{xcolor}
\RequirePackage{enumitem}
\RequirePackage[most]{tcolorbox}
\RequirePackage{tabularx,booktabs,array}
\RequirePackage{fancyhdr}
\RequirePackage{needspace}
\RequirePackage{geometry}
\RequirePackage[hidelinks]{hyperref}

\geometry{a4paper,top=17mm,bottom=18mm,left=23mm,right=23mm}
\definecolor{notegray}{gray}{0.39}
\setlength{\parindent}{0pt}
\setlength{\parskip}{2.5pt plus 1pt minus 0.5pt}
\setlist[itemize]{topsep=2pt,itemsep=1pt,parsep=0pt,leftmargin=1.45em}
\setlist[enumerate]{topsep=2pt,itemsep=2pt,parsep=0pt,leftmargin=1.65em}
\setcounter{secnumdepth}{0}
\renewcommand{\arraystretch}{1.12}

% Zahlavi: vlevo rocnik | globalni tema, vpravo cislo: nazev hodiny
% Zapati zustava zcela prazdne.
\pagestyle{fancy}
\fancyhf{}
\renewcommand{\headrulewidth}{0.35pt}
\renewcommand{\footrulewidth}{0pt}
\setlength{\headheight}{14pt}
\newcommand{\nastavPripravu}[4]{%
  \fancyhead[L]{\small #1\enspace\textbar\enspace #2}%
  \fancyhead[R]{\small #3: #4}%
}

% Minimalisticke cernobile boxy. Kazda definice zvlast.
\tcbset{commonlesson/.style={enhanced,breakable=false,colback=white,colframe=black,
  colbacktitle=white,coltitle=black,boxrule=0.45pt,
  arc=0.8mm,outer arc=0.8mm,left=3mm,right=3mm,top=1.6mm,bottom=1.6mm,
  toptitle=1mm,bottomtitle=0.7mm,fonttitle=\bfseries,
  attach boxed title to top left={xshift=2.5mm,yshift=-1.15mm},
  boxed title style={colback=white,colframe=white,boxrule=0pt,left=1mm,right=1mm}}}
\newtcolorbox{zadaniBox}[1][]{commonlesson,title={Zadání},#1}
\newtcolorbox{prikladBox}[1][]{commonlesson,title={Příklad},#1}
\newtcolorbox{definiceBox}[1][]{commonlesson,title={Definice},#1}

% Spolecny box pro zadani a jeho reseni, bez mezery mezi castmi.
% Pouziti:
%   \begin{prikladReseniBox}
%     Zadani ulohy.
%     \tcblower
%     \begin{align*} ... \end{align*}
%   \end{prikladReseniBox}
% Spodni cast ma jemny sedy podklad a delici linku; lze v ni sazet
% vice radku matematiky. Obe casti maji jen jeden spolecny ramecek.
\tcbset{solutionlesson/.style={
  commonlesson,
  enhanced jigsaw,
  breakable,
  colbacklower=black!3,
  segmentation style={solid,draw=black!25,line width=0.35pt},
  fontlower=\small\color{notegray},
  before lower={\noindent\textbf{Řešení:}\enspace},
  leftlower=3mm,rightlower=3mm
}}
\newtcolorbox{prikladReseniBox}[1][]{solutionlesson,title={Příklad},#1}
\newtcolorbox{zadaniReseniBox}[1][]{solutionlesson,title={Zadání},#1}

% Pomucky pro autora pripravy: sedive poznamky; \reseni pro starsi pripravy
\newcommand{\poznamka}[1]{\par\begingroup\small\color{notegray}\itshape
  \textnormal{\textbf{Poznámka pro učitele.}}\enspace #1\par\endgroup}
\newcommand{\reseni}[1]{\par\begingroup\small\color{notegray}
  \textbf{Očekávané řešení:}\enspace #1\par\endgroup}
\newcommand{\blok}[1]{\needspace{5\baselineskip}\subsection*{#1}}
\newcommand{\vyrok}[1]{\uv{#1}}

\usepackage{tikz}
\usetikzlibrary{arrows.meta}
\nastavPripravu{1. ročník}{Funkce}{60}{Definiční obor a obor hodnot}
\begin{document}

\poznamka{\textbf{Jádro na 45 minut:} návrat k algebraickým výrazům
3 min, definice a tabulka upozornění 10 min, první dvě úlohy
10 min, obor hodnot z tabulky a předpisu 9 min, čtení grafu 7 min,
samostatné procvičení a kontrola 6 min. Doplňkové úlohy jsou rezerva. V této hodině neuč
obecné určování oboru hodnot složitých racionálních funkcí;
čtení složitějších grafů přijde v hodině 61.}

\blok{Co už známe z algebraických výrazů?}
U lomených výrazů jsme určovali \textbf{podmínky přípustných hodnot
proměnné}. Například výraz
\[
\frac{x^2-4}{x-2}=\frac{(x-2)(x+2)}{x-2}\qquad (x\ne2)
\]
má smysl jen při $x\ne2$. Po krácení získáme $x+2$, ale původně
zakázaná hodnota $x=2$ zůstává vyloučena. Stejnou otázku si nyní
klademe u funkce: \emph{pro které argumenty je zadána a jakých hodnot
může skutečně nabývat?}

Připomeň si také příklad z minulé hodiny: obsah čtverce je
$S(a)=a^2$. Algebraicky umíme vypočítat $a^2$ pro každé $a\in\mathbb R$,
ale \textbf{délka strany čtverce musí být kladná}, tedy v tomto modelu
$a>0$. Předpis sám proto nemusí vystihovat všechna omezení úlohy.

\blok{Definiční obor: které vstupy jsou dovoleny?}
\begin{definiceBox}
\textbf{Definiční obor} funkce $f$, značený $D(f)$, je množina všech
argumentů, pro které je funkce zadána. U funkce $f\colon A\to B$
platí $D(f)=A$.
\end{definiceBox}

Je-li funkce zadána jen předpisem bez dalšího omezení a pracujeme-li
v reálných číslech, hledáme obvykle \emph{přirozený definiční obor
předpisu} — největší podmnožinu $\mathbb R$, na níž má výraz smysl.
Pokud zadání stanoví menší množinu vstupů, řídíme se \textbf{zadáním},
ne pouze vzorcem.

\medskip
\textbf{Na co si dát pozor při určování $D(f)$?}

\begingroup
\small
\renewcommand{\arraystretch}{1.28}
\begin{tabularx}{\linewidth}{@{}>{\raggedright\arraybackslash}p{.205\linewidth}>{\raggedright\arraybackslash}p{.345\linewidth}X@{}}
\toprule
\textbf{Typ výrazu} & \textbf{Co musí platit} & \textbf{Příklad}\\
\midrule
Zlomek $\dfrac{P(x)}{Q(x)}$
& Jmenovatel nesmí být nulový: $Q(x)\ne0$.
& $\dfrac1{x-2}$: $x\ne2$.\\
Druhá odmocnina $\sqrt{P(x)}$
& Odmocňovaný výraz musí být nezáporný: $P(x)\ge0$.
& $\sqrt{x+2}$: $x\ge-2$.\\
Třetí odmocnina $\sqrt[3]{P(x)}$
& Záporná čísla jsou dovolena. Samotná třetí odmocnina obor neomezuje.
& $\sqrt[3]{-8}=-2$; $\sqrt[3]{x-5}$ pro každé $x\in\mathbb R$.\\
Sinus, kosinus $\sin x$, $\cos x$
& Argument může být libovolné reálné číslo. \emph{Hodnoty} obou funkcí jsou mezi $-1$ a $1$.
& $\sin 2$ má smysl; $\sin x>1$ nemá řešení.\\
\emph{Pro zajímavost:} $\arcsin P(x)$, $\arccos P(x)$
& Zde naopak musí být \emph{vstup} v $\langle-1;1\rangle$:
$-1\le P(x)\le1$.
& $\arcsin(2x)$: $-\frac12\le x\le\frac12$.\\
\bottomrule
\end{tabularx}
\endgroup

\emph{Pozor na záměnu:} u $\sin$ a $\cos$ omezení $\langle-1;1\rangle$
patří \textbf{oboru hodnot}, nikoli definičnímu oboru. Setká-li se
v jednom předpisu více podmínek, musí platit \textbf{současně}:
např. $f(x)=\dfrac{\sqrt{x-1}}{x-3}$ vyžaduje $x\ge1$ a $x\ne3$.


\needspace{12\baselineskip}
\begin{prikladReseniBox}[title={Příklad 1: připomenutí podmínek}]
Urči přirozené definiční obory funkcí
$f(x)=\dfrac{x+1}{x-3}$ a $g(x)=\sqrt{x+2}$.
\tcblower
U zlomku nesmí být jmenovatel nulový:
\[
x-3\ne0\quad\Longrightarrow\quad
D(f)=\mathbb R\setminus\{3\}.
\]
U druhé odmocniny musí být odmocňovaný výraz nezáporný:
\[
x+2\ge0\quad\Longrightarrow\quad
D(g)=\langle-2;+\infty).
\]
V obou případech používáme pravidla známá z úprav algebraických výrazů.
\end{prikladReseniBox}

\begin{prikladReseniBox}[title={Příklad 2: zadání je důležitější než samotný vzorec}]
Je dána funkce $h(x)=x^2$ pro $x\in\langle-1;2\rangle$.
Urči její definiční obor. Změnil by se, kdybychom žádný obor neuvedli?
\tcblower
Zadaná funkce má $D(h)=\langle-1;2\rangle$. Výraz $x^2$
sám o sobě je definovaný na celém $\mathbb R$, ale zde je funkce
záměrně omezena. Bez dalšího omezení bychom uvažovali
$D(h)=\mathbb R$.
\end{prikladReseniBox}

\poznamka{Při příkladu 2 naváž na 59. hodinu: stejný předpis na dvou
různých definičních oborech určuje různé funkce. Rozlišuj také
\uv{mohu výraz spočítat} od \uv{patří vstup do zadaného modelu}.}

\blok{Obor hodnot: které výstupy skutečně nastanou?}
\begin{definiceBox}
\textbf{Obor hodnot} funkce $f$, značený $H(f)$, je množina všech
funkčních hodnot, kterých funkce skutečně nabývá:
\[
H(f)=\{f(x)\mid x\in D(f)\}.
\]
Je-li $f\colon A\to B$, potom $H(f)\subseteq B$; cílová množina $B$
se nemusí celá vyčerpat.
\end{definiceBox}

\begin{prikladReseniBox}[title={Příklad 3: konečný definiční obor}]
Funkce $p\colon\{-2,0,1,3\}\to\{0,1,4,9,16\}$ je dána
$p(x)=x^2$. Urči $D(p)$, $H(p)$ a cílovou množinu.
\tcblower
Dosadíme všechny přípustné argumenty:
\[
\begin{array}{c|rrrr}
x&-2&0&1&3\\ \hline
p(x)&4&0&1&9
\end{array}
\]
Proto $D(p)=\{-2,0,1,3\}$, $H(p)=\{0,1,4,9\}$ a cílová množina
je $\{0,1,4,9,16\}$. Číslo $16$ je v cílové množině, ale funkce
ho na zadaném oboru nenabývá.
\end{prikladReseniBox}

\begin{prikladReseniBox}[title={Příklad 4: obor hodnot ze známého předpisu}]
Je dáno $q(x)=x^2$ pro $x\in\langle-2;1\rangle$.
Urči definiční obor a obor hodnot.
\tcblower
$D(q)=\langle-2;1\rangle$. Druhá mocnina je nezáporná;
nejmenší hodnotu $q(0)=0$ a největší hodnotu $q(-2)=4$ funkce
nabývá. Protože pro každé $y\in\langle0;4\rangle$ můžeme zvolit
$x=-\sqrt y\in\langle-2;0\rangle$, nastávají i všechny mezilehlé
hodnoty. Tedy $H(q)=\langle0;4\rangle$.
\end{prikladReseniBox}

\blok{Čtení definičního oboru a oboru hodnot z grafu}
V grafu hledáme přípustné \textbf{vodorovné souřadnice} (osa $x$) a
skutečně nabývané \textbf{svislé souřadnice} (osa $y$).
Plný krajní bod do grafu patří; prázdný krajní bod nikoli.

\begin{center}
\begin{tikzpicture}[x=.98cm,y=.73cm,>=Latex]
\draw[step=1,gray!20,very thin] (-3,-2) grid (3,4);
\draw[->] (-3.25,0)--(3.35,0) node[right] {$x$};
\draw[->] (0,-2.25)--(0,4.32) node[above] {$y$};
\foreach \x in {-3,-2,-1,1,2,3} {\draw (\x,.08)--(\x,-.08) node[below=2pt,font=\scriptsize] {$\x$};}
\foreach \y in {-2,-1,1,2,3,4} {\draw (.08,\y)--(-.08,\y) node[left=2pt,font=\scriptsize] {$\y$};}
\draw[very thick] (-2,-1)--(2,3);
\fill (-2,-1) circle (2.9pt);
\draw[thick,fill=white] (2,3) circle (3.2pt);
\node[below left,font=\scriptsize] at (-2,-1) {$(-2;-1)$};
\node[above right,font=\scriptsize] at (2,3) {$(2;3)$};
\end{tikzpicture}
\end{center}

\begin{prikladReseniBox}[title={Příklad 5: krajní body grafu}]
U vyznačeného úseku grafu $y=x+1$ urči $D(f)$, $H(f)$ a rozhodni,
zda funkce nabývá hodnoty $3$.
\tcblower
Vodorovné souřadnice začínají číslem $-2$ včetně, ale k $2$ se
graf pouze blíží. Proto $D(f)=\langle-2;2)$. Svislé souřadnice
začínají hodnotou $-1$ včetně a končí před hodnotou $3$:
$H(f)=\langle-1;3)$. Hodnoty $3$ funkce nenabývá.
\end{prikladReseniBox}

\poznamka{Zde záměrně nevyvozuj monotonii ani průsečíky s osami;
patří následujícím hodinám. Nech nejprve jednoho žáka ukázat
\emph{na ose $x$} definiční obor a jiného \emph{na ose $y$} obor hodnot.}

\blok{Samostatné procvičení}
\begin{zadaniReseniBox}[title={Samostatná práce}]
\begin{enumerate}[label=\alph*)]
\item Urči přirozený definiční obor $a(x)=\dfrac{1}{x+4}$.
\item Urči přirozený definiční obor $b(x)=\sqrt{5-x}$.
\item Pro $c(x)=2x-1$ a $D(c)=\{-1,0,3\}$ urči $H(c)$.
\item Pro $d(x)=x^2$ na $\langle0;3\rangle$ urči $H(d)$.
\end{enumerate}
\tcblower
\begin{enumerate}[label=\alph*)]
\item $D(a)=\mathbb R\setminus\{-4\}$, protože $x+4\ne0$.
\item $D(b)=(-\infty;5\rangle$, protože $5-x\ge0$.
\item $c(-1)=-3$, $c(0)=-1$, $c(3)=5$; $H(c)=\{-3,-1,5\}$.
\item $H(d)=\langle0;9\rangle$: na nezáporném intervalu
od $0$ do $3$ druhá mocnina nabývá všech hodnot od $0$ do $9$.
\end{enumerate}
\end{zadaniReseniBox}

\blok{Volitelná rezerva}
\begin{zadaniReseniBox}[title={Rozlišení předpisu a skutečného modelu}]
Cena výroby $n$ kusů je $C(n)=80n+250$ Kč, kde $n\in\mathbb N_0$.
Urči $D(C)$ a zapiš $H(C)$. Proč nemá smysl dosazovat $n=1{,}5$?
\tcblower
$D(C)=\mathbb N_0$ a $H(C)=\{250+80n\mid n\in\mathbb N_0\}$.
Počet hotových kusů je celé nezáporné číslo; $n=1{,}5$
není přípustným argumentem tohoto modelu.
\end{zadaniReseniBox}

\blok{Domácí procvičení}
\begin{zadaniReseniBox}[title={Domácí úkol}]
\begin{enumerate}[label=\alph*)]
\item Urči přirozený definiční obor $f(x)=\dfrac{x^2-1}{x-1}$.
Můžeš po zkrácení dosadit $x=1$?
\item Urči přirozený definiční obor $g(x)=\sqrt{2x-6}$.
\item Funkce $h(x)=3-x$ je zadána na $D(h)=\langle0;4\rangle$.
Urči $H(h)$.
\item Funkce $k(x)=x^2$ má $D(k)=\{-3,-1,0,2\}$.
Urči $H(k)$.
\item Funkce $m(x)=x^2$ je zadána na $\langle-3;1\rangle$.
Urči $H(m)$ a zdůvodni krajní hodnoty.
\item Pro funkci $u(x)=4-x$ na $D(u)=\langle0;5)$
urči $H(u)$. Která z krajních hodnot je do něj zahrnuta?
\end{enumerate}
\tcblower
\begin{enumerate}[label=\alph*)]
\item $D(f)=\mathbb R\setminus\{1\}$. Po krácení získáme $x+1$,
ale $x=1$ zůstává vyloučeno.
\item $2x-6\ge0$, tedy $D(g)=\langle3;+\infty)$.
\item $H(h)=\langle-1;3\rangle$; hodnoty $3$ a $-1$ nastanou
v krajních bodech intervalu.
\item $k(-3)=9$, $k(-1)=1$, $k(0)=0$, $k(2)=4$;
$H(k)=\{0,1,4,9\}$.
\item $H(m)=\langle0;9\rangle$: minimum $m(0)=0$ a maximum
$m(-3)=9$; všechny hodnoty mezi nimi funkce nabývá.
\item $H(u)=(-1;4\rangle$. Číslo $4$ je zahrnuto ($u(0)=4$),
číslo $-1$ nikoli (odpovídalo by nepřípustnému $x=5$).
\end{enumerate}
\end{zadaniReseniBox}

\end{document}
